% commentaire.tex — f165r, vv. 208–251 (conversion initiale des travaux antérieurs)
\comm{224}{οὗτος σὺ ποῖ θεῖς;}{formule identique au vers 854 des \italique{Guêpes}.}
\comm{224}{εἰς τὸ τῶν σεμνῶν θεῶν·}{Le scholiaste :
\citComm{ἀντὶ τοῦ `εἰς <τὸ> τῶν ᾿Ερινύων'. ἄσυλον γὰρ εἶχον οἱ καταλαμβάνοντες τὰ ἱερὰ τῶν Ἐρινύων.}}
\comm{232}{ψιλὸς αὖ στρατεύσομαι.}{ψιλός, comme \italique{levis} en latin, signifie à la fois \italique{lisse}, \italique{sans poil}, et \italique{léger}, \italique{armé à la légère}. (Note de Willems.)\\
Cette note sur ψιλός me sera l'occasion de digresser sur l'un des plus beaux échanges à bâtons rompus de la Tragédie grecque : chez Sophocle (\italique{Ajax}, v. 1120 sq.), entre Ménélas et Teucer :
\citComm{Μενέλαος — Ὁ τοξότης ἔοικεν οὐ σμικρὸν φρονεῖν.\\	
Τεύκρος — Οὐ γὰρ βάναυσον τὴν τέχνην ἐκτησάμην.\\
Μενέλαος — Μέγ᾽ ἄν τι κομπάσειας, ἀσπίδ᾽ εἰ λάβοις.\\
Τεύκρος — Κἂν \textbf{ψιλὸς} ἀρκέσαιμι σοί γ᾽ ὡπλισμένῳ.\\
Μενέλαος — Ἡ γλῶσσά σου τὸν θυμὸν ὡς δεινὸν τρέφει.\\
Τεύκρος — Ξὺν τῷ δικαίῳ γὰρ μέγ᾽ ἔξεστιν φρονεῖν.}Ménélas s'étonne de ce que Teucer se permette de lui adresser des reproches – à lui le roi de Sparte ! –, vil archer de son état --- Mais l'archerie est loin d'être méprisable lui rétorque Teucer --- Ménélas ironise : Que serait-ce s'il portait un bouclier --- La réplique de Teucer fuse : « Même \italique{ψιλὸς}, je tiendrais bon devant toi et ton [royal] armement ». Comprendre : même \italique{fantassin de l'armée légère} ; c'est-à-dire l'échelon le plus bas chez les militaires – encore plus bas qu'archer –, poste souvent réservé à des esclaves et donc le plus opposé à la royale position de Ménélas et le plus inégal en terme d'armements. C'en est trop pour Ménélas qui s'indigne de cette langue par trop pendue. C'est alors que Teucer lance la fameuse et cinglante réplique : « La fierté est permise, quand on a pour soi la justice ! »}
\comm{235}{Κλεισθένην.}{Clisthène, l'efféminé notoire, personnage qui apparaîtra dans la pièce au vers 571 -- et sera d'ailleurs reconnue comme femme -- avant de prêter allégeance au genre féminin. 
Cette tête de Turc d'Aristophane est raillée comme : une pseudo-femme (\italique{Nuées} 355, \italique{Oiseaux} 829-831, \italique{Thesmophories} 571) ; en tout cas femme en amour (\italique{Lysistrata} 621, 1091-1092, \italique{Thesmophories} 574, \italique{Grenouilles} 48) ; imberbe comme femme (\italique{Acharniens} 117-121, \italique{Grenouilles} 422-424) ; Clisthène ne contrebalance même pas ces innombrables défauts – qui ne sont pas tous du ressort de sa volonté – par une misogynie affectée : il trouve encore le moyen d'être traître à son sexe en s'alliant au genre féminin (\italique{Lysistrata} 621, \italique{Thesmophories} 575-576, 602).}
\comm{237}{δελφάκιον}{est un diminutif d'ἡ δέλφαξ, le porc \italique{mature}, opposé au χοῖρος, \italique{le porcelet}.
N'en déplaise à certains commentateurs, il n'est fait aucune allusion ici au sexe féminin -- pas plus qu'au vers 1061 de \italique{Lysistrata}. Et d'abord, ce serait absurde : il n'est question que du côté pile ; côté épilé par flambage. 
L'image est ailleurs : c'est celle de viande de cochon dont on retire les poils par le feu.
Cf. la \italique{Souda} (δ 205):\citComm{δελφάκιον : ὁ μικρὸς χοῖρος. Ἀριστοφάνης: οἴμοι κακοδαίμων, δελφάκιον γέγονα. τοῦτο εἴρηκεν ἀφευθεὶς τὸν πρωκτόν. ἀντὶ τοῦ φλογισθείς. μετὰ γὰρ τὸ τυθῆναι τὰ δελφάκια φλογίζονται, ἵνα ψιλωθῶσιν.}(On pourra aussi se reporter à Homère\nc{\italique{Iliade}, IX, 467-8 et \italique{Odyssée}, II, 300 ; XIV, 73-5} ou Eschyle\nc{Fragment 310 de l'édition de S. Radt, \italique{Tragicorum Graecorum Fragmenta}, vol. 3 (1985).}.)}
\comm{239}{τὴν κέρκον φυλάττου νυν ἄκραν.}{Euripide reprend l'image du δελφάκιον flambé ; il parle bien-sûr de l'\italique{olisbos} de Mnésiloque.}
\comm{241–2}{οἴμοι τάλας. ὕδωρ ὕδωρ ὦ γείτονες / πρὶν ἀντιλαβέσθαι πρωκτὸν ἓτερον τῆς φλογός.}{Mnésiloque est littéralement en train de prendre feu ; toute sa réplique est un pastiche d'un appel au secours, comme si sa propre maison s'embrasait\nc{La scholie au vers 242 : (πρωκτὸν·) ὡς εἰ ἔλεγεν ``οἰκίαν''.} et que celles, contiguës, de ses voisins allaient prendre feu à leur tour (sauf qu'il s'agit ici non d'οἶκος mais de πρωκτός)... il appelle donc à puiser de l'eau « ὕδωρ ! ὕδωρ ! » Horace prodigue ce conseil dans ses \italique{Épîtres} (1, 18, 84) :
\citComm{\latin{Tua res agitur, paries cum proximus ardet}\\
\francais{C'est ton intérêt qui est en jeu, lorsque la maison de ton voisin brûle\nc{J'emprunte la traduction au magistral \italique{Dictionnaire des sentences latines et grecques} de Renzo Tosi (traduit de l'italien par Rebecca Lenoir). La consultation de ce dictionnaire me fut d'une grande aide pour découvrir, en ce passage comme en d'autres, des ponts entre la langue d'Aristophane et l'univers des proverbes.}.}}
L'image de l'incendie qui se propage de maison en maison était si fréquente dans le monde latin et roman\nc{Voir le \italique{Dictionnaire} sus-cité (§524).} qu'on a peine à croire qu'elle ne le fût pas aussi dans le monde grec -- témoin ce passage d'Aristophane.}
\comm{243}{θάρρει — τί θαρρῶ;}{Plaute, \italique{Curculio}, IV, 3, 21 : \latin{bellator vale — quid valeam ?}}
\comm{247}{μὴ φροντίσῃς· ἕτερος γὰρ αὐτὰ σπογγιεῖ.}{Les Grecs se servaient d'éponges pour la toilette et, dans un usage détourné, pouvaient même les utiliser en guise de torche-cul comme cela ressort d'une scène des \italique{Grenouilles} (vv. 481 sqq.). Qui est cet ἕτερος ? Réponse : le voisin d'à côté : « La situation est bien celle qu'on attend : en cas d'incendie, les voisins, soucieux de leur propre intérêt, commençaient par faire la chaîne pour apporter de l'eau. Puis, le feu une fois éteint, il se trouvait presque toujours un voisin complaisant pour aider le propriétaire (ou le locataire) malchanceux à nettoyer sa maison et à y remettre de l'ordre. » (Taillardat\nc{Jean Taillardat, «~Aristophane, \italique{Thesmophories} v. 247, 289, 721 sqq. : essai d'interprétation du texte transmis~», in \italique{Revue des Études Grecques}, t. 85, fasc. 406-408, juillet-décembre 1972, p. XVII.
Le grand Taillardat avait donc compris le sens de ce passage avant même que le \nieme{242} vers, lacunaire, ne fût rétabli par la découverte d'un papyrus.}.)}
\comm{248}{οἰμώξετ᾽ ἄρ᾽ εἴ τις τὸν ἐμὸν πρωκτὸν πλυνεῖ.}{πλύνειν signifie au propre \italique{nettoyer} (et fait donc écho au σπογγιεῖ du vers précédent) et par suite \italique{dire du mal} (\italique{passer un savon} en français ; l'on manque toujours de la délicatesse requise lorsqu'on s'occupe de la toilette d'autrui).
Le sel de la boutade de Mnésiloque nous échappe... 
Mnésiloque souhaite-t-il bien du courage à qui voudrait \italique{s'attaquer} -- au prime sens de πλύνειν : \italique{commencer de nettoyer} -- à son fondement ? Ingrate tâche qui serait aussitôt à recommencer. Travail de Sisyphe qu'on pourrait rapprocher du proverbe grec πλίνθον πλύνειν, « laver une brique »\nc{Voir Photius ο 360 (\italique{s.v.} ὄνου πόκοι), ο 363 (ὄνου πόκαι) et l'apparat conséquent des \latin{testimonia} et autres \latin{loci paralleli} de l'édition Theodoridis. Le proverbe fut répertorié entre autres par Plutarque, \italique{Corpus paroemiographorum Graecorum}, I, περὶ τῶν ἀδυνάτων ; Diogenianus (VII, 50), Zenobius (VI, 48). Le scholiaste des \italique{Guêpes} (280b) le cite : <λίθον ἕψεις·> ὡς καὶ τὸ `\textbf{πλίνθον πλύνεις}' καὶ `χύτραν ποικίλλεις' καὶ `εἰς ὕδωρ γράφεις' καὶ `Αἰθίοπα λευκαίνεις' καὶ `κατὰ θαλάττης σπείρεις' καὶ τὰ τοιαῦτα. C'est d'ailleurs en écho à ce proverbe qu'Aristophane rapproche au vers 514 du \italique{Ploutos} les verbes πλινθουργεῖν et πλύνειν.}. Dans le même genre d'idée, Aristophane emploie au vers 604 des \italique{Guêpes} l'expression qui paraît proverbiale πρωκτὸς λουτροῦ περιγιγνόμενος, « le c.. l'emporte sur l'eau lustrale »\nc{Bdélycléon lançait cette sentence à la barbe de Philocléon, saoul de discours et qui, quelque effort qu'il fît pour soigner sa glossolalite, n'en retomberait pas moins un jour dans son travers.
Le proverbe devait donc peu ou prou signifier qu'un naturel grossier, le πρωκτός, retournait inévitablement à sa fange ; en français : un \italique{indécrottable}. Il est repris tel quel par les parémiographes Diogenianus (III, 53) et Michel Apostolius (XIV, 78), Hésychius (π 4144), Photius (π 1450), Suidas (π 2951). La source principale doit être d'anciens commentaires au vers 604 des \italique{Guêpes}. (Cf. les scholies aux manuscrits d'Aristophane ; l'on y apprend que Καλλίστρατος lui-même avait signalé l'expression comme proverbiale.)}.}
